\documentclass[10pt,a4paper]{article}
\usepackage[utf8]{inputenc}
\usepackage[T1]{fontenc}
\usepackage[spanish,es-nodecimaldot]{babel}
\usepackage{lmodern}
\usepackage[margin=1.85cm,headheight=14pt]{geometry}
\usepackage{xcolor}
\usepackage{booktabs,tabularx,array,longtable}
\usepackage{enumitem}
\usepackage{fancyhdr}
\usepackage{hyperref}
\usepackage{listings}
\usepackage{microtype}
\definecolor{navy}{HTML}{173F67}
\definecolor{blue}{HTML}{276FA8}
\definecolor{soft}{HTML}{F3F6F9}
\definecolor{line}{HTML}{D5DFE8}
\hypersetup{colorlinks=true,linkcolor=navy,urlcolor=blue}
\pagestyle{fancy}
\fancyhf{}
\lhead{\small Seminario de Programación - Periodo 2}
\rhead{\small Instituto Jorge Robledo}
\cfoot{\small \thepage}
\setlength{\parindent}{0pt}
\setlength{\parskip}{5pt}
\setlist[itemize]{leftmargin=5mm,itemsep=2pt,topsep=3pt}
\setlist[enumerate]{leftmargin=6mm,itemsep=3pt,topsep=3pt}
\newcommand{\ruleline}{\par\noindent\color{line}\rule{\linewidth}{0.6pt}\color{black}\par}
\newcommand{\statusbox}[2]{\noindent\fcolorbox{line}{soft}{\parbox{0.94\linewidth}{\textbf{#1}\\#2}}}
\newcommand{\scorebox}[2]{\fcolorbox{line}{soft}{\parbox{0.43\linewidth}{\centering\textcolor{navy}{\Large\textbf{#1}}\\[-1mm]\small #2}}}
\lstset{basicstyle=\ttfamily\small,backgroundcolor=\color{soft},frame=single,rulecolor=\color{line},breaklines=true,columns=fullflexible,showstringspaces=false}
\begin{document}
\begin{center}
{\color{navy}\Large\bfseries Guía docente previa a la clase}\\[2mm]
{\small Seminario de Programación - 29 de julio de 2026}
\end{center}
\ruleline

\section*{Objetivo de la sesión}
Revisar el estado real de cada proyecto ya subido, confirmar equipos o parejas, ejecutar el último SHA, acordar una meta verificable y dejar un commit nuevo con evidencia individual.
\section*{Estado antes de la clase}
\scriptsize
\begin{longtable}{p{3.0cm}p{2.6cm}p{1.5cm}p{1.5cm}p{5.2cm}}
\toprule
\textbf{Estudiante} & \textbf{Proyecto} & \textbf{Repo} & \textbf{Estado} & \textbf{Primera prioridad}\\
\midrule
Juan Pablo Arango Giraldo & Sistema de pedidos Delivery & Vinculado & 2.11/5 & Crear y conectar una hoja CSS \\
Jerónimo Mazo López & Sistema de registro estudiantil & Pendiente & No evaluable & Vincular repositorio y subir versión mínima. \\
Jerónimo Rodríguez Peña & Carrito de restaurante & Vinculado & 3.55/5 & Implementar validaciones explícitas de entrada \\
Pedro Pablo Arbeláez Escobar & Cajero automático didáctico & Vinculado & 1.40/5 & Crear y conectar una hoja CSS \\
Samuel Chavarriaga Avendaño & Formulario multipágina con persistencia & Pendiente & No evaluable & Vincular repositorio y subir versión mínima. \\
Pablo Jaramillo Palacio & Registro de actividades & Vinculado & 2.09/5 & Crear y conectar una hoja CSS \\
Alejandro Rico Páramo & Gestor de tareas académicas & Pendiente & No evaluable & Vincular repositorio y subir versión mínima. \\
Tomás González Giraldo & Sistema de inventario escolar & Pendiente & No evaluable & Vincular repositorio y subir versión mínima. \\
Alejandro Rincón Torres & Calculadora y panel de notas & Pendiente & No evaluable & Vincular repositorio y subir versión mínima. \\
\bottomrule
\end{longtable}
\normalsize
\section*{Orden recomendado}
\begin{enumerate}
\item Proyectos con error sintáctico o que no ejecutan.
\item Proyectos con repositorio pero baja integración visible.
\item Proyectos compartidos: confirmar integrantes, roles y autoría.
\item Proyecto de referencia: cerrar faltantes y profundizar CSS.
\item Estudiantes sin repositorio: vincular cuenta, pareja y versión mínima.
\end{enumerate}

\newpage
\section*{Protocolo por estudiante o pareja (10-12 minutos)}
\begin{enumerate}
\item Confirmar nombres, grupo, proyecto, repositorio y modalidad individual o grupal.
\item Ejecutar la página principal del último SHA copiado.
\item Verificar flujo principal, error visible y persistencia si aplica.
\item Pedir explicación: entrada, función, estado, DOM y salida.
\item Seleccionar una meta concreta para hoy.
\item Realizar una modificación en vivo y un commit descriptivo.
\item Registrar resultado, bloqueador, integrantes y siguiente paso.
\end{enumerate}
\section*{Reglas para parejas}
\begin{itemize}
\item Máximo dos estudiantes por proyecto durante este periodo.
\item Un repositorio compartido puede representar a ambos.
\item README obligatorio con integrantes y roles.
\item La nota técnica del repositorio puede ser común; explicación, autoría y modificación en vivo son individuales.
\item Cada integrante debe aportar commits propios o demostrar claramente su parte.
\end{itemize}
\section*{Referencia técnica actual}
Carrito de restaurante de Jerónimo Rodríguez Peña: 3.55/5.00 y 71\%. Se usa como comparación por la mayor evidencia técnica, no porque el proyecto esté terminado.
\section*{Producto mínimo al terminar la clase}
\begin{itemize}
\item Equipo o modalidad individual confirmada.
\item Repositorio vinculado y página principal ejecutada.
\item Meta concreta registrada.
\item Cambio verificable y commit nuevo.
\item Guía individual entregada o compartida.
\end{itemize}
\end{document}
